\documentclass{article}
\usepackage[T1]{fontenc}
\usepackage{lmodern}
\usepackage{graphicx}
\usepackage{amsmath,amssymb}
\usepackage[a4paper,margin=2cm]{geometry}
\usepackage[absolute,overlay]{textpos}
\setlength{\TPHorizModule}{1mm}
\setlength{\TPVertModule}{1mm}
\usepackage[indonesian]{babel}
\usepackage{hyperref}
\setlength{\emergencystretch}{2em}
% unit: Halaman 40

\begin{document}

% === Halaman 40 (llm) ===

\section*{7. Cipher}

\begin{itemize}
    \item Algoritma untuk enkripsi dan dekripsi pesan
    \item Aturan (\textit{rule}) untuk \textit{enchipering} dan \textit{dechipering}, atau berupa fungsi matematika yang digunakan untuk enkripsi dan dekripsi pesan.
\end{itemize}

Contoh:
\begin{itemize}
    \item Enkripsi: Geser tiga huruf ke kanan \quad (rule)
    \item Dekripsi: Geser tiga huruf ke kiri \quad (rule)
    \item \[ E(p) = (p + k) \pmod{26} \quad(\text{fungsi matematika}) \]
    \item \[ D(c) = (c - k) \pmod{26} \quad(\text{fungsi matematika}) \]
\end{itemize}

\begin{itemize}
    \item \textcolor{red}{Classical cipher}: Caesar cipher, Vigenere cipher, Playfair Cipher, Enigma cipher, dll
    \item \textcolor{red}{Modern cipher}: DES, AES, Blowfish, Serpent, RSA, ElGamal, RC4, RC5, A5, dll
\end{itemize}

\end{document}
